%% Manju Vallayil — CV for ICE Clinical Data Mapping Analyst (Graduate)
%% Compile with: pdflatex ice_cv.tex  (run twice)

\documentclass[10.5pt,a4paper]{article}

% ── Packages ──────────────────────────────────────────────────────────────────
\usepackage[T1]{fontenc}
\usepackage[utf8]{inputenc}
\usepackage[left=1.8cm, right=1.8cm, top=1.5cm, bottom=1.8cm]{geometry}
\usepackage{charter}
\usepackage{microtype}
\usepackage{parskip}
\usepackage{enumitem}
\usepackage{xcolor}
\usepackage{titlesec}
\usepackage{array}
\usepackage{tabularx}
\usepackage[most]{tcolorbox}
\usepackage[
  colorlinks=true,
  urlcolor=accent,
  linkcolor=accent,
  citecolor=accent
]{hyperref}
\usepackage{fontawesome5}

% ── Colours ───────────────────────────────────────────────────────────────────
\definecolor{accent}{HTML}{3E0097}
\definecolor{accentlt}{HTML}{EDE8F7}
\definecolor{mid}{HTML}{444444}
\definecolor{light}{HTML}{999999}

% ── Section headings ──────────────────────────────────────────────────────────
\titleformat{\section}
  {\color{accent}\small\bfseries\scshape}{}{0em}{\MakeUppercase}
  [\vspace{1pt}\color{accent}\rule{\linewidth}{1.5pt}]
\titlespacing{\section}{0pt}{16pt}{7pt}

% ── List style ────────────────────────────────────────────────────────────────
\setlist[itemize]{
  leftmargin=1.3em, itemsep=2pt, topsep=3pt, parsep=0pt,
  label={\color{accent!60}\small\textbullet}
}
\newcommand{\tabitem}{{\color{accent!70}\textbullet}\enspace}

% ── Experience left-border box ────────────────────────────────────────────────
\newenvironment{expbox}{%
  \begin{tcolorbox}[
    blanker, breakable,
    left=10pt,
    borderline west={2.5pt}{0pt}{accentlt},
    before upper={\parskip=2pt},
    top=0pt, bottom=0pt
  ]
}{%
  \end{tcolorbox}\vspace{4pt}
}

% ── Role / Edu macros ─────────────────────────────────────────────────────────
\newcommand{\role}[4]{%
  {\bfseries #1}\hfill{\small\color{accent}#2}\\[-1pt]
  {\small\color{mid}#3}\hfill{\small\color{light}#4}\\[3pt]%
}

\newcommand{\edu}[4]{%
  {\bfseries #1}\hfill{\small\color{accent}#2}\\[-1pt]
  {\small\color{mid}#3}\hfill{\small\color{light}#4}\\[2pt]%
}

\pagestyle{empty}
\setlength{\parindent}{0pt}

% ══════════════════════════════════════════════════════════════════════════════
\begin{document}

% ── Full-bleed header ─────────────────────────────────────────────────────────
\vspace*{-1.5cm}
\begin{tcolorbox}[
  enhanced, colback=accent, colframe=accent,
  arc=0pt, boxrule=0pt,
  left=1.8cm, right=1.8cm, top=16pt, bottom=14pt,
  enlarge left by=-1.8cm, enlarge right by=-1.8cm,
  width=\textwidth+3.6cm
]
  \begin{minipage}[t]{0.54\linewidth}
    {\color{white}\LARGE\bfseries Manju Vallayil}\\[5pt]
    {\color{white!65!accent}\small
      Data Analyst\enspace·\enspace
      Clinical Informatics \& ML\enspace·\enspace
      PhD, Computer Science}
  \end{minipage}%
  \hfill
  \begin{minipage}[t]{0.43\linewidth}
    \raggedleft\small\color{white!80!accent}
    \faEnvelope\enspace
      \href{mailto:manjuvallayil@gmail.com}{\textcolor{white!80!accent}{manjuvallayil@gmail.com}}\\[3pt]
    \faLinkedin\enspace
      \href{https://linkedin.com/in/manju-vallayil-phd-5b1b7a142/}{\textcolor{white!80!accent}{linkedin/manju-vallayil-phd}}\\[3pt]
    \faGlobe\enspace
      \href{https://manju.vallayil.com}{\textcolor{white!80!accent}{manju.vallayil.com}}\\[3pt]
    \faMapMarker*\enspace Auckland, New Zealand
  \end{minipage}
\end{tcolorbox}

\vspace{12pt}

% ── Profile ───────────────────────────────────────────────────────────────────
\section{Profile}

Data analyst and AI researcher with a PhD in Computer Science and hands-on experience
working with clinical datasets. Applied ML to EEG, clinical, and psychological data in a
tinnitus detection project, handling preprocessing, quality validation, and cross-source
documentation. An earlier career in software quality assurance built a solid grounding in
data verification, test documentation, and systematic accuracy checks. Broader research in
NLP and explainable AI has consistently required close attention to detail and clear
communication with non-technical stakeholders. Bringing this combined background in clinical
data, analytical rigour, and structured documentation to a healthcare informatics role.

% ── Key Skills ────────────────────────────────────────────────────────────────
\section{Key Skills}

\renewcommand{\arraystretch}{1.35}
\begin{tabular}{@{}p{0.48\linewidth}p{0.48\linewidth}@{}}
  \textbf{\color{mid}Data Analysis \& Quality}
    & \textbf{\color{mid}Research \& Methods} \\
  \tabitem Data validation, cleaning \& quality assurance
    & \tabitem Literature review \& evidence synthesis \\
  \tabitem Clinical data preprocessing \& feature engineering
    & \tabitem Qualitative \& mixed-methods research \\
  \tabitem SQL-based data verification
    & \tabitem Technical writing \& documentation \\
  \tabitem Data mapping, classification \& annotation
    & \tabitem Research ethics \& compliance \\[20pt]
  \textbf{\color{mid}Tools \& Languages}
    & \textbf{\color{mid}Domain Knowledge} \\
  \tabitem Python (data analysis, ML)
    & \tabitem Clinical ML (tinnitus, EEG datasets) \\
  \tabitem SQL, Microsoft Excel
    & \tabitem AI for healthcare \& education \\
  \tabitem scikit-learn, spaCy, Hugging Face
    & \tabitem Stakeholder communication \& reporting \\
  \tabitem LaTeX, Git
    & \tabitem NLP \& information extraction \\
\end{tabular}
\renewcommand{\arraystretch}{1}

% ── Experience ────────────────────────────────────────────────────────────────
\section{Experience}

\begin{expbox}
\role{Research Assistant — AI Modelling for Tinnitus Detection}
     {Mar 2022 – Mar 2023}
     {Auckland University of Technology}
     {Auckland, NZ}
\begin{itemize}
  \item Applied machine learning to EEG, clinical, and psychological datasets to model
        factors influencing tinnitus detection and treatment outcomes in an interdisciplinary
        research setting.
  \item Performed data preprocessing, feature engineering, and model evaluation —
        identifying patterns across heterogeneous clinical data sources and resolving
        inconsistencies between datasets.
  \item Maintained clear documentation of data processing decisions and model outputs,
        communicating findings to clinical stakeholders without prior ML background.
  \item Built and documented a minimum viable prototype (MVP) demonstrating practical
        application of AI-driven research findings to a clinical audience.
\end{itemize}
\end{expbox}

\begin{expbox}
\role{Research Assistant — AI for Education \& Community Innovation}
     {Mar 2025 – Dec 2025}
     {Auckland University of Technology}
     {Auckland, NZ}
\begin{itemize}
  \item Conducted data analysis on adoption patterns across the \textbf{TechTahi} project,
        working with mixed qualitative and quantitative datasets to identify meaningful trends.
  \item Designed and implemented research methodology, applying rigorous quality checks to
        ensure consistency and accuracy of analytical outputs.
  \item Co-authored peer-reviewed publication in \textit{Education Sciences} (May 2026),
        contributing to evidence synthesis and interpretation of outcomes.
\end{itemize}
\end{expbox}

\begin{expbox}
\role{Research Officer — AI-Driven Language Learning Technology}
     {Sep 2021 – Dec 2024}
     {Auckland University of Technology}
     {Auckland, NZ}
\begin{itemize}
  \item Co-developed \textbf{Te Kōhanga o te Tūī}, an AI-enabled bilingual learning platform
        combining speech recognition and NLP for early childhood acquisition of
        te reo Māori and English.
  \item Contributed to NLP research, technical architecture documentation, and evaluation
        of model performance within a live deployment context.
  \item Liaised across language specialists, developers, and community stakeholders to
        translate research outputs into functional, accessible technology.
\end{itemize}
\end{expbox}

\begin{expbox}
\role{Business Analyst}
     {Oct 2017 – Mar 2019}
     {Tureya Limited}
     {Auckland, NZ}
\begin{itemize}
  \item Led requirements gathering and stakeholder communication for mobile application
        development — producing user needs documentation, data flow diagrams, and UX
        wireframes supporting MVP delivery.
  \item Developed and executed test plans, verifying system logic and data flows against
        specified requirements before client handover.
\end{itemize}
\end{expbox}

\begin{expbox}
\role{Software Quality Assurance Engineer}
     {May 2012 – Feb 2014}
     {TechWyse Internet Marketing}
     {Cochin, India}
\begin{itemize}
  \item Built data-driven test suites and performed SQL-based backend validation to verify
        data integrity and system accuracy across web applications.
\end{itemize}
\end{expbox}

\begin{expbox}
\role{Junior Software Quality Assurance Engineer}
     {Nov 2010 – Apr 2012}
     {ART Technology \& Software India Pvt.\ Ltd}
     {Cochin, India}
\begin{itemize}
  \item SQL-based backend data validation, test case documentation, and system testing of
        e-commerce platforms integrated with Microsoft RMS.
\end{itemize}
\end{expbox}

% ── Selected Publications ─────────────────────────────────────────────────────
\section{Selected Publications}

\begin{itemize}

  \item \textbf{Vallayil, M.}, Nand, P., Yan, W., et al.
        ``CARAG: A Context-Aware Retrieval Framework for Fact Verification, Integrating
        Local and Global Perspectives of Explainable AI.''
        \textit{Applied Sciences}, Feb 2025.
        \href{https://www.mdpi.com/2076-3417/15/4/1970}{\small[link]}

  \item \textbf{Vallayil, M.}, Nand, P., Yan, W., Allende-Cid, H.
        ``Explainability of Automated Fact Verification Systems: A Comprehensive Review.''
        \textit{Applied Sciences}, Nov 2023.
        \href{https://www.mdpi.com/2076-3417/13/23/12608}{\small[link]}

  \item \textbf{Vallayil, M.}, Nand, P., Yan, W.
        ``Explainable AI through Thematic Clustering and Contextual Visualization:
        Advancing Macro-Level Explainability in AFV Systems.''
        \textit{ACIS}, Dec 2024.
        \href{https://aisel.aisnet.org/acis2024/101/}{\small[link]}

\end{itemize}

{\small\color{light}Full list (10+ papers, datasets, technical reports):
\href{https://scholar.google.com/citations?user=6xGoEsMAAAAJ}{Google Scholar}\enspace·\enspace
\href{https://www.researchgate.net/profile/Manju-Vallayil-2}{ResearchGate}}

% ── Education ─────────────────────────────────────────────────────────────────
\section{Education}

\edu{Doctor of Philosophy (PhD) — Computer Science}
    {Mar 2020 – Aug 2025}
    {Auckland University of Technology (AUT)}
    {Auckland, NZ}
{\small\textit{Thesis:} \href{https://openrepository.aut.ac.nz/items/33de7437-821c-4a19-8539-0d38c6848d04}{Advancing Explainable AI: A Global Context-Aware Evidence Retrieval Framework for Interpretable Fact Verification}}

\vspace{6pt}

\edu{Master of Information Technology}
    {Mar 2016 – Apr 2019}
    {Eastern Institute of Technology}
    {New Zealand}

\edu{Bachelor of Technology (BTech) — Computer Science \& Engineering}
    {Mar 2004 – Dec 2009}
    {Government Engineering College, Palakkad}
    {India}

% ── Awards & Recognition ──────────────────────────────────────────────────────
\section{Awards \& Recognition}

\renewcommand{\arraystretch}{1.4}
\begin{tabular}{@{}p{0.78\linewidth}r@{}}
  \href{https://verify.skilljar.com/c/ngmh5b4cn9um}{Claude Code in Action Certification — Anthropic}
    & \textcolor{accent}{\small\bfseries 2026} \\
  AUT Doctoral Journal Writing Grant
    & \textcolor{accent}{\small\bfseries 2025} \\
  AUT DCT Faculty Postgraduate Summer Research Award {\small(×2)}
    & \textcolor{accent}{\small\bfseries 2023, 2024} \\
  Futures Pacific (CIIRID: IDEA) — AUT-PUCV Chile Scholarship {\small(×2)}
    & \textcolor{accent}{\small\bfseries 2021, 2023} \\
  AUT ECMS Research Scholarship
    & \textcolor{accent}{\small\bfseries 2022} \\
  PhD Full Tuition Fee Waiver
    & \textcolor{accent}{\small\bfseries 2021--2024} \\
\end{tabular}

\end{document}
